\par\noindent\textbf{(1)} 令 \(u=1/r\)，中心吸引力 \(F=-k/r^2\) 的Binet式为 \(u''+u=mk/L^2\)，解为 \(u=(mk/L^2)[1+e\cos(\phi-\phi_0)]\)。利用能量积分得 \(e^2=1+2EL^2/(mk^2)=1+(mv_\infty^2b/k)^2>1\)，因此轨道为双曲线。
\par\noindent\textbf{(2)} 取近心点在 \(\phi=\phi_0\)，则 \(r_p=L^2/[mk(1+e)]\)。又记 \(h=mv_\infty^2b/k\)，代入可得 \(r_p=(k/(mv_\infty^2))[\sqrt{1+h^2}-1]\)，两式相等并且 \(r_p\to0\) 当 \(b\to0\)；碰撞中心在点势下出现奇异。
\par\noindent\textbf{(3)} 无穷远处 \(u\to0\) 要求 \(\cos(\phi_\infty-\phi_0)=-1/e\)。两条渐近线之间的速度偏转角为 \(\chi=2\arcsin(1/e)=2\arctan[k/(mv_\infty^2b)]\)。等价地 \(b=[k/(mv_\infty^2)]\cot(\chi/2)\)，高能或大 \(b\) 时散射角很小。
\par\noindent\textbf{(4)} 轴对称散射有 \(d\sigma/d\Omega=(b/\sin\chi)|db/d\chi|\)。代入 \(b=C\cot(\chi/2)\)、\(C=k/(mv_\infty^2)\)，得 \(\frac{d\sigma}{d\Omega}=\frac{C^2}{4\sin^4(\chi/2)}=\frac{k^2}{4m^2v_\infty^4\sin^4(\chi/2)}\)。总截面在小角因长程势而发散，需要实验角截断。
